% Quarantined v4 source-correct strict-open dependency for P06.2.
% The legacy fibre-slicing proof has been deleted, rather than patched in
% place.  The active polycylinder proof cites the principal v4 proposition,
% whose local moment step uses CLD Proposition 15.1.1(b) in the forward
% direction.  These new exact bytes are not production-verified.

\begin{proof}[Proof of \cref{prop:p062-arbitrary-field-annulus-v4-source-correct}]
For \(s\in(r,R)\), let \(\eta_s\) be the \(s\)-Gauss point. Testing the
current equation on compactly supported functions of
\(-\log|T|\), and using the weight-one split-torus integration formula,
shows that
\begin{equation}\label{eq:p062-annulus-radial-affine-v4-source-correct}
             u\longmapsto h(\eta_{\mathrm e^{-u}})
\end{equation}
has zero distributional second derivative on
\((-\log R,-\log r)\). Thus it is affine. This is exactly the
one-dimensional part of the split-torus calculation used in
\cref{prop:p062-rational-centred-traces}; unlike the disc case, there is
no endpoint at \(T=0\) and its affine slope is not forced to vanish.

We next prove that every rigid point has the value prescribed by
\eqref{eq:p062-annulus-radial-affine-v4-source-correct}. Let \(y\) be rigid, put
\(E=\mathscr H(y)\), \(s=|T(y)|\), and let
\[
                 P(T)=T^d+\sum_{i=0}^{d-1}c_iT^i
\]
be the monic irreducible polynomial of \(T(y)\). Completeness of \(K\)
gives a unique extension of its absolute value to every finite extension.
Every conjugate of \(T(y)\), counted with inseparable multiplicity, has
absolute value \(s\). Consequently
\begin{equation}\label{eq:p062-annulus-coefficient-bounds-v4-source-correct}
             |c_i|\leq s^{d-i},
             \qquad |c_0|=s^d=:S\in|K^\times|.
\end{equation}
Strict dominance of the constant term for \(|T|<s\) and of the leading
term for \(|T|>s\) gives
\begin{equation}\label{eq:p062-annulus-polynomial-tube-v4-source-correct}
 X_y:=\{|P(T)|<S\}\subseteq\{|T|=s\}\subset A^-_{r,R}.
\end{equation}
The polynomial \(P\) induces a finite flat morphism of degree \(d\)
\[
             p\colon X_y\longrightarrow D^-_S,
             \qquad Z=P(T).
\]
Let \(H=\operatorname{Tr}_p(h|_{X_y})\), with analytic local degrees in
the fibre trace. By \cref{prop:p062-finite-flat-continuous-trace}, \(H\)
is continuous and \(\mathrm d'\mathrm d''[H]=0\). The arbitrary-field
strict-disc theorem \cref{prop:p062-arbitrary-field-disc} makes \(H\)
constant. The zero fibre is the single analytic point \(y\), with degree
\(d\), so
\begin{equation}\label{eq:p062-annulus-trace-zero-v4-source-correct}
                              H=d\,h(y).
\end{equation}

Let \(\xi_S\) denote the boundary Gauss point of the target closed disc.
The graded Gauss calculation used in the proof of
\cref{prop:p062-arbitrary-field-disc}, now at \(t=s\), gives
\begin{equation}\label{eq:p062-annulus-single-gauss-fibre-v4-source-correct}
                      P^{-1}(\xi_S)=\{\eta_s\},
                      \qquad \deg_{\eta_s}(P)=d.
\end{equation}
Indeed, the initial polynomial
\(\operatorname{in}_s(P)(U)\) has degree \(d\); the rational-function
extension generated by \(U\) over that generated by
\(\operatorname{in}_s(P)(U)\) has degree \(d\), including its inseparable
part. The graded fundamental inequality gives local degree at least \(d\),
and finite flatness leaves neither degree nor another point in the fibre.

Choose a neighbourhood \(W\) of \(\eta_s\) whose closure is contained in
the annulus. Since \(P\) is finite and
\eqref{eq:p062-annulus-single-gauss-fibre-v4-source-correct} is a singleton fibre,
there is an analytic neighbourhood \(V\) of \(\xi_S\) such that
\(P^{-1}(V)\subset W\). Applying the continuity part of
\cref{prop:p062-finite-flat-continuous-trace} to
\(P^{-1}(V)\to V\) shows, as \(q\uparrow S\), that
\[
             H(\xi_q)\longrightarrow d\,h(\eta_s).
\]
Together with \eqref{eq:p062-annulus-trace-zero-v4-source-correct}, this yields
\begin{equation}\label{eq:p062-annulus-rigid-to-skeleton-v4-source-correct}
                              h(y)=h(\eta_s).
\end{equation}
Rigid points are dense locally on the strict annulus. The affinity of
\eqref{eq:p062-annulus-radial-affine-v4-source-correct},
\eqref{eq:p062-annulus-rigid-to-skeleton-v4-source-correct}, and continuity give
\eqref{eq:p062-annulus-affine-v4-source-correct} at every analytic point.
Uniqueness is immediate because \(\log|T|\) is nonconstant on the
skeleton.
\end{proof}

\begin{proof}[Proof of \cref{thm:p062-arbitrary-field-polycylinder-v4-source-correct}]
We first work with the open polycylinder. Apply
\cref{prop:p062-strict-product-fibre-slicing-v4-research} to any coordinate
projection. On a disc fibre, \cref{prop:p062-arbitrary-field-disc} shows
that the restriction of \(h\) is constant. Hence \(h\) is independent of
every disc coordinate.
If \(q=0\), this already proves that \(h\) is constant, and there is
nothing more to show.

Fixing the \(K\)-rational origin in every disc factor and replacing the
disc coordinates successively shows, by the preceding constancy, that
\(h\) is the pullback of a continuous function on the product of the
annulus factors. We continue to denote that function by \(h\).

On an annulus fibre, \cref{prop:p062-arbitrary-field-annulus-v4-source-correct}
shows that the restriction of \(h\) is affine in the logarithmic radius.
In particular, it has at an arbitrary fibre point the same value as at the
Gauss point with the same radius. The slicing proposition is valid over
every base point, with its possibly nonrigid complete residue field. It
remains to verify that successive replacements in distinct coordinates
have a common, order-independent endpoint.

Fix radii \(\mathbf s=(s_1,\ldots,s_q)\) in the open radius rectangle. For
a Laurent polynomial
\(
 f=\sum_{\nu\in\mathbf Z^q}a_\nu S^\nu
\), successive one-variable Gauss extensions, in any order, give
\begin{equation}\label{eq:p062-iterated-gauss-v2}
                  |f|_{\mathbf s}
                   =\max_{\nu}|a_\nu|
                         \prod_{j=1}^q s_j^{\nu_j}.
\end{equation}
Indeed, each step takes the maximum over one exponent and the finitely
many maximum operations commute. At an intermediate step the coefficients
may lie in the completed residue field of the coordinates not yet
replaced. The one-variable Gauss norm is still defined by the same maximum
formula, and passing to that completion does not change the result.

The formula extends from Laurent polynomials to every Laurent series
converging on a fixed compact subannulus. Such a series is the Banach-norm
limit of finite Laurent sums, and every bounded multiplicative seminorm is
continuous for that norm. It follows that every order of replacement
induces the same seminorm on each compact subannulus, hence the same
analytic point on the open polyannulus. By
\eqref{eq:p062-iterated-gauss-v2}, that point is the canonical
multivariate Gauss point \(\sigma(\mathbf s)\).

The section \(\mathbf s\mapsto\sigma(\mathbf s)\) is continuous. Indeed,
on a compact logarithmic subrectangle, the norm of a Laurent polynomial is
the maximum of finitely many continuous monomials in \(\mathbf s\), and
the preceding Banach approximation gives the norm of every analytic
function as a locally uniform limit of such functions. This is precisely
continuity for the Berkovich topology.

Now let \(z\) be an arbitrary point of the open polyannulus and put
\(s_j=|S_j(z)|\). Replacing the coordinates one at a time does not change
the value of \(h\), while the preceding calculation identifies the final
point independently of order. Therefore
\[
                     h(z)=h\bigl(\sigma(\mathbf s)\bigr).
\]
It follows that \(h\) factors through the standard tropicalization:
\begin{equation}\label{eq:p062-polycylinder-factor-v4-source-correct}
 h=\phi\circ F,
 \qquad
 F=(-\log|S_1|,\ldots,-\log|S_q|),
\end{equation}
where
\[
             \phi(\mathbf u)
                =h\bigl(\sigma(\mathrm e^{-u_1},\ldots,
                                      \mathrm e^{-u_q})\bigr).
\]
The continuity just proved for the Gauss section and continuity of \(h\)
show that \(\phi\) is continuous on the open rectangle
\(
 \prod_j(-\log R'_j,-\log r_j).
\)

If \(p>0\), choose a compactly supported smooth positive top-degree form
\(\gamma\) in the product of the disc factors with
\(\int\gamma>0\). If \(p=0\), put \(\gamma=1\) and
\(\int\gamma=1\). For an arbitrary compactly supported Lagerberg form
\(\beta\) of bidegree \((q-1,q-1)\) on the logarithmic rectangle, use as
a test form the product of \(\gamma\) with the pullback of \(\beta\).
Since \(\gamma\) has top degree in the disc variables,
\(\mathrm d'\gamma=\mathrm d''\gamma=0\). The weight-one split-torus
integration formula and Fubini therefore convert
\(
 \mathrm d'\mathrm d''[h]=0
\)
into
\(
 \mathrm d'\mathrm d''[\phi]=0
\)
in the distributional Lagerberg sense on that rectangle; the nonzero
factor \(\int\gamma\) cancels. Equivalently,
every distribution
\(
 \partial^2\phi/\partial u_i\partial u_j
\)
vanishes. Convolution on relatively compact Euclidean balls makes every
regularization affine; compatibility on overlaps and continuity show that
\(\phi\) itself is affine. This proves
\eqref{eq:p062-polycylinder-affine-v4-source-correct}.

Each \(\log|S_j|\) is the logarithm of the norm of an analytic unit.
Thus the displayed function is ordinary CLD-smooth with zero Hessian,
weakly smooth and Gubler--Jell--Rabinoff harmonic, and classically
pluriharmonic. Uniqueness follows by restriction to the standard
polyannulus skeleton.
\end{proof}
